\section*{الفصل \ref{ch:b2:scalar-light-model} --- النموذج السُّلَّمي للضوء}

\begin{solution}{exo:b2:scalar-light-model:1}
$(AB) = \qty{1.5}{cm}$؛ و $25\,000$ طول موجة؛ و \qty{0.5}{cm} أكثر من الهواء، أي $8300$
طول موجة، و $\Delta\varphi = 2\pi \times 8300$، وتأخير \qty{17}{ps}؛ ويكون التأخير طول
موجة واحدة عند $e = \lambda/(n - 1) = \qty{1.2}{\micro m}$.
\end{solution}

\begin{solution}{exo:b2:scalar-light-model:2}
$\ell_c = \lambda^2/\Delta\lambda$ أو $c/\Delta\nu$، و $\tau_c = \ell_c/c$: فالأبيض \qty{0.9}{\micro m} (\qty{3}{fs})؛
والصوديوم \qty{1.7}{cm} (\qty{60}{ps})؛ والثنائي الباعث \qty{16}{\micro m} (\qty{50}{fs})؛ والثنائي الليزري \qty{3}{mm}
(\qty{10}{ps})؛ والمستقر \qty{300}{km} (\qty{1}{ms}). فعند \qty{1}{mm}: مصباح الصوديوم،
والثنائي الليزري، والليزر؛ وعند \qty{10}{m}: الليزر المستقر وحده.
\end{solution}

\begin{solution}{exo:b2:scalar-light-model:3}
$10^{-3}/3.06 \times 10^{-19} = \qty{3.3e15}{}$ في الثانية. وأما المصباح: $I = 5/4\pi \times 9 =
\qty{44}{mW/m^2}$، والبؤبؤ \qty{28}{mm^2}: \qty{1.2}{\micro W}، أي $3.5 \times 10^{12}$ فوتونًا في الثانية،
و $3.5 \times 10^{11}$ في \qty{0.1}{s}، والتقلب $2 \times 10^{-6}$: أي أملس تمامًا.
\end{solution}

\begin{solution}{exo:b2:scalar-light-model:4}
المضبوط $\sqrt{1 + 10^{-4}} - 1 = \qty{49.9988}{\micro m}$، والمحوري القريب \qty{50}{\micro m}: أي
مئة طول موجة؛ \omterm{def:b2:scalar-light-model:path}{وجبهة الموجة} كرة نصف قطرها \qty{1}{m}؛ ويبلغ
الحدّ الرباعي $r^4/8d^3$ المقدارَ $\lambda$ عند $r = (8\lambda d^3)^{1/4} = \qty{4.5}{cm}$.
\end{solution}

\begin{solution}{exo:b2:scalar-light-model:5}
(أ) للكرة التي نصف قطرها $f$ ومركزها $F'$، على مستوي العدسة،
الطور $-\pi r^2/\lambda f$ بالنسبة إلى المحور (محوريًّا قريبًا): وهو طور
العدسة بالضبط. (ب) والطور الكلي $\pi r^2/\lambda\,(1/d - 1/f) = -\pi r^2/\lambda d'$: أي موجة
متقاربة عند $d'$ حيث $1/d' = 1/f - 1/d$. (ج) ويضيف الزجاج $(n - 1)e(r)
\times2\pi/\lambda$: $e(r) = e_0 - r^2/2(n - 1)f$؛ والسهم $10^{-4}/0.1 = \qty{1}{mm}$. (د) والطور
المعرَّف بترديد $2\pi$ هو الطور نفسه --- عند طول موجة
التصميم؛ وعند غيره لا تعود الدرجات أطوال موجة صحيحة.
\end{solution}

\begin{solution}{exo:b2:scalar-light-model:6}
(أ) من \omterm{def:b2:scalar-light-model:path}{جبهة موجة} مستوية عند $x = -D$ إلى $(x, y)$ ثم إلى $F(f, 0)$: $D + x
+ \sqrt{(x - f)^2 + y^2} = D + f$، ومن ثَمّ $\sqrt{(x - f)^2 + y^2} = f - x$، أي $y^2 = 4fx$،
و $F$ في البؤرة. (ب) والدائرة $x = 2f - \sqrt{4f^2 - y^2} \approx y^2/4f +
y^4/64f^3$ توافقه حتى المرتبة الثانية فقط. (ج) و $y^4/64f^3 = \qty{1.6}{\micro m}$، مضاعفًا
بالانعكاس: \qty{3}{\micro m}، أي ستة أطوال موجة --- ومن هنا المرايا
المكافئة. (د) وعند $\lambda \sim \qty{1}{cm}$ يكون الخطأ نفسه جزءًا من ألف من
طول الموجة.
\end{solution}

\begin{solution}{exo:b2:scalar-light-model:7}
(أ) $L(x) = n_1\sqrt{a^2 + x^2} + n_2\sqrt{b^2 + (d - x)^2}$: $L' = n_1\sin\theta_1 -
n_2\sin\theta_2 = 0$. (ب) وبجوار المسار المستقر يكون للمسارات المجاورة
أطوار متساوية فتتجمع أمواجها؛ وفي غير ذلك تتلاشى. (ج) و $n_1 = n_2$:
$\sin\theta_1 = \sin\theta_r$. (د) \omterm{ex:b2:scalar-light-model:lens}{والتطابق البؤري} يعني أن لكل الأشعة من $A$ إلى $A'$
مسيرات متساوية: أي أسرة كاملة من المسارات المستقرة، لا مسارًا واحدًا.
\end{solution}

\begin{solution}{exo:b2:scalar-light-model:8}
(أ) $I = 2I_0(1 + \cos\Delta\varphi)$: $4I_0$ و $0$. (ب) و $N = 10^7$: فالمتوسط $2I_0$، والباقي
$\sim 2I_0/\sqrt N = 6 \times 10^{-4}$ منه. (ج) وخلال \qty{10}{ps} يكون الطور
مجمّدًا: فيرى هدبًا تقفز كل \qty{0.1}{ns}. (د) وبليزرات
$\tau_c \sim \qty{1}{\micro s}$ وكاشف أسرع من ذلك: فالهدب موجودة
مدة ميكروثانية وقد صُوّرت.
\end{solution}

\begin{solution}{exo:b2:scalar-light-model:9}
(أ) عند الزوايا الصغيرة: $\tan i \approx n\tan r$، فتبدو الأشعة آتية من $h/n$.
(ب) \qty{0.90}{m}؛ \omterm{def:b2:scalar-light-model:path}{والمسير الضوئي} $n \times 1.2 = \qty{1.6}{m}$، أي \qty{0.4}{m} أكثر.
(ج) و $n \times 1 = \qty{1.33}{m}$. (د) وعند الزوايا الكبيرة يشتدّ الانكسار
ويصغر العمق الظاهري.
\end{solution}

\begin{solution}{exo:b2:scalar-light-model:10}
(أ) موجتان متساويتان بفرق طور $2\pi\nu\delta/c$. (ب) والمتوسط
على النطاق:
\[
\frac1{\Delta\nu}\int\cos\frac{2\pi\nu\delta}c\,\dd\nu = \cos\frac{2\pi\nu_0\delta}c\,\operatorname{sinc}\frac{\pi\Delta\nu\,\delta}c .
\]
(ج)
$I_{\max, \min} = 2I_0(1 \pm |\operatorname{sinc}|)$: $V = |\operatorname{sinc}(\pi\Delta\nu\delta/c)|$،
وهو معدوم عند $\delta = c/\Delta\nu$. (د) و $\Delta\nu = c\Delta\lambda/\lambda^2 = \qty{1.9e13}{Hz}$: $\delta =
\qty{16}{\micro m}$.
\end{solution}

\begin{solution}{exo:b2:scalar-light-model:11}
(أ) $100$ فوتون في الثانية على \qty{38}{mm^2}: \qty{2.6e6}{m^{-2}.s^{-1}}، $\times4 \times
10^{-19}$ جول: \qty{1e-12}{W/m^2}. (ب) و $10^{-13.1} \times 1000 = \qty{8e-11}{W/m^2}$، أي $2 \times 10^8$
فوتونًا في المتر المربع والثانية، و $8000$ في الثانية إلى البؤبؤ
--- أي فوق (أ) بكثير: فحدّ العين الحقيقي تحدده الخلفية ويحدده
التوزّع على عصيّات كثيرة. (ج) و $(200/7)^2 = 800$، أي $+7.3$ من القدر: $m \approx
13$. (د) و $N$ فوتونًا تتقلب بمقدار $\sqrt N$: فتنمو نسبة الإشارة إلى الضجيج
كالجذر التربيعي لزمن التعريض.
\end{solution}

\begin{solution}{exo:b2:scalar-light-model:12}
(أ) $d(1 + r^2/2d^2 - r^4/8d^4 + \dots) - d$. (ب) $r^4/8d^3 < \lambda/4$. (ج) $r <
(2\lambda d^3)^{1/4}$: \qty{3.2}{cm} من أجل \qty{1}{m}؛ و \qty{1}{mm} من أجل \qty{1}{cm}، بزاوية
\qty{0.1}{rad}. (د) والزيغان الكروي هو هذا الحدّ الرباعي؛ وأما
الزيغانات الأخرى فهي أبناء عمه خارج المحور.
\end{solution}

\begin{solution}{pb:b2:scalar-light-model:1}
\textbf{1.} $\nu_1 = \qty{5.0934e14}{Hz}$، و $\nu_2 = \qty{5.0882e14}{Hz}$: $\Delta\nu = \qty{5.2e11}{Hz}$.

\textbf{2.} $c/\Delta\nu_{\text{طبيعي}} = \qty{30}{m}$.

\textbf{3.} $\sqrt{k_BT/m} = \qty{425}{m/s}$: $\Delta\nu_D \approx 2\nu v/c \approx \qty{1.4}{GHz}$، أي
مئة ضعف العرض الطبيعي.

\textbf{4.} $c/\Delta\nu_D \approx \qty{20}{cm}$؛ ومع التصادمات، نحو \qty{2}{cm}.

\textbf{5.} $\lambda^2/\Delta\lambda = 589^2/0.6$ نانومتر $= \qty{0.58}{mm}$: أي من التوافق إلى
التعاكس كل \qty{0.29}{mm}؛ ونحو $1000$ هدبة في الضربة.

\textbf{6.} يتذبذب التباين بدور \qty{0.58}{mm} في فرق
المسير (وينعدم عند $0.29$، و $0.87$، و \qty{1.45}{mm}، \dots) داخل غلاف
يخبو على بضعة سنتيمترات: أي نحو خمسين ضربة، ثم لا شيء.

\textbf{7.} \qty{3e18}{} فوتون في الثانية؛ و \qty{2.1}{eV}، وهو أصفر
الطيف.

\textbf{8.} لا: فالذرات مستقلة، والأطوار عشوائية. وخطا المصباح
الواحد يأتيان أيضًا من ذرات مختلفة ولا يتداخل أحدهما مع
الآخر: فالضربات نظاما هدب يتجمعان في \emph{الشدة}.

\textbf{9.} $\lambda^2/\Delta\lambda = \qty{0.85}{mm}$؛ و $c/\Delta\nu = \qty{300}{m}$.

\textbf{10.} $\qty{3.3e15}{}$ في الثانية، متباعدة $c/N = \qty{90}{nm}$ على امتداد الحزمة:
أي عشرة آلاف فوتون في طول ترابط واحد. فالتداخل ليس
لقاء فوتونات: بل كل فوتون يتداخل مع نفسه.

\textbf{11.} $\Delta\nu = c\Delta\lambda/\lambda^2 = \qty{360}{GHz}$؛ \omterm{def:b2:maxwell-equations:operators}{وتباعد} الأنماط $c/2nL = \qty{43}{GHz}$:
أي نحو ثمانية أنماط.

\textbf{12.} $e \times 3.3 \times 10^{15} = \qty{0.5}{mA}$.

\textbf{13.} \qty{2}{mm} $>$ \qty{0.85}{mm}: فلا هدب؛ و \qty{10}{m} $<$ \qty{300}{m}: فهدب.

\textbf{14.} النبضة التي مدتها نانوثانية طولها \qty{30}{cm}، أي أطول بكثير من
\omterm{prop:b2:scalar-light-model:coherence}{طول الترابط} \qty{0.85}{mm}: فعرض الخط، لا النبضة، هو الذي يحدد
الترابط؛ ولا تتداخل النبضات المتعاقبة إلا إذا نجا طور الليزر
من فترة الإطفاء --- وهو لا ينجو.

\textbf{15.} الضوء المترابط المتبعثر على الجدار الخشن يتداخل عند
الشبكية بفروق مسير عشوائية: وهو الترقيط. وترابط المصباح
على سُلَّم الميكرومتر يتوسط ذلك كله فيمحوه.

\textbf{16.} $\lambda^2/\Delta\lambda = \qty{17}{\micro m}$: أي أقصر بكثير من أن يعطي ترقيطًا على
جدار خشن --- فلا ترقيط، وصورة أنعم.

\textbf{17.} $(n - 1)e = \qty{1}{mm}$: أي $1540$ هدبة.

\textbf{18.} $r = \qty{6.65}{\degree}$؛ والمسير الزائد $\approx e(n/\cos r - \cos(i - r)/\cos r)
= 3.021 - 2.010 = \qty{1.011}{mm}$، أي \qty{11}{\micro m} أكثر من غير المائلة: أي نحو
عشرين هدبة أخرى.

\textbf{19.} خمسة بكسلات في الهدبة؛ وعند بكسل واحد في الهدبة يكامل البكسل
دورًا كاملًا فينهار التباين.

\textbf{20.} \qty{1e-12}{W} على البكسل $\times$ \qty{10}{ms} $= \qty{1e-14}{J}$: أي $3 \times 10^4$
فوتون، والضجيج $0.6\%$.

\textbf{21.} $I = 2I_0[1 + \cos(2\pi\delta/\lambda)]$: أي هدبة واحدة في كل طول موجة من
فرق المسير.

\textbf{22.} تتوسط العين الحركة الكيلوهرتزية إلى حقل
منتظم؛ أما الثنائي الضوئي فيتابعها.

\textbf{23.} $(n - 1) \times 10 = \qty{2.9}{mm}$: أي $4500$ هدبة؛ وتغيّر الضغط بنسبة $1\%$
يزيحها $45$ هدبة --- فمقياس التداخل مقياس ضغط ما لم
يُفرَّغ.

\textbf{24.} معدل الفوتونات هو نفسه تقريبًا (\qty{2.1}{eV} مقابل \qty{1.9}{eV})؛
والهدب تتعلق بالسعات والأطوار، وأما عدّ الفوتونات فلا يدخل إلا في
الضجيج.

\textbf{25.} المصباح: سنتيمترات، يحدّها التعريض الدوبلري
والتصادمي، والهدب حتى نحو \qty{1}{cm}؛ والمؤشر الرخيص: \qty{0.85}{mm}،
تحدّه أنماطه المتعددة؛ والليزر المستقر: \qty{300}{m}، يحدّه
عرض خطه.
\end{solution}
